\typeout{NT FILE chap4.tex}%

\chapter{Solution Proposal}
\label{cha:solution_proposal}

\prependtographicspath{{5-Figures/Covers/}}

% epigraph configuration
\epigraphfontsize{\small\itshape}
\setlength\epigraphwidth{12.5cm}
\setlength\epigraphrule{0pt}

\section{Proposed System Architecture}
\label{sec:sol_architecture}
\ntindex[Template]{}
    {
        The proposed platform follows a three-layer architecture that separates user interface, application logic and data persistence, while remaining compatible with the technologies and requirements established in earlier chapters.
        At a high level, this structure is designed to support incremental evolution from a prototype to a long-lived institutional system, without forcing disruptive changes to the overall desgn.
    }

    {
        At the presentation layer, a React-based web application delivers role-specific dashboards for students, company representatives, professors, coordinators, and the secretariat.
        The UI will be composed of reusable components for common views, which simplifies maintaining consistent behavior and visual language across roles.
        For student authentication, the platform will integrate with the university's identity provider using OAuth, so that students
        can log in with their institutional accounts and their profile and enrollment information can be imported without manual registration.
    
    }

    {
        The application layer exposes a REST API that encapsulates all business logic. 
        During the initial prototyping phase, this layer can rely on managed services (such as Supabase) for basic authentication
        and CRUD operations, so that core workflows can be validated quickly with real users.
        As the system stabilizes, a Spring Boot-based backed assumes responsibility for the domain logic.
    }

    {
        The data layer is implemented with PostgreSQL, using a relational schema such as the one described in Section 4.3.
        Referential integrity is enforced with foreign keys to prevent inconsistent states, and \gls{RLS} is used to constrain which rows each role can access.
        External services complement this layer: a transactional email provider is used for reliable delivery and tracking of notifications, and an object storage
        service holds upload documents such as CVs, protocols and reports. 
    }

    {
        Together, these elements allow AI-assisted prototypes to be refined into a secure, maintainable platform rather than being discarded.
    }

    \begin{figure}[htbp]
            \centering
            \includegraphics[height=5in]{MyFigures/architecture_diagram}
            \caption{High-level three layer architecture diagram of the platform}
            \label{fig:Figure_Architecture_Diagram}
        \end{figure}


\section{Process Modelling (BPMN)}
\label{sec:sol_bpmn}
\ntindex[Template]{}
    {
        To complement the textual use case model and clarify the most critical workflows, three BPMN diagrams were
        created and included in the appendices. These models focus on the processes that generate most of the coordination
        effort in the system: company registration, proposals submission and validation, and student seriation and assginements.
        They are intentionally simple and somewhat "raw", serving as a first iteration that captures main actors, decisions,
        and state transitions without yet modelling every possible exception or alternate path.
    }

    {
        The first diagram describes the registration of a company representative and the creation or association
        of a company entity in the platform. 
        It shows how an employee accesses an external registration page, submits personal details, and either selects
        and existing company, or creates a new one, while the system validates input, stores employee and company data,
        and sends verification e-mails.
    }

    {
        The second diagram models the proposal submission and validation workflow for companies and professors.
        It begins when the submitter enters an open edition, fills in the proposal form, and decides whether to save a draft
        or submit; the system then saves the proposal and notifies coordinators that a new item is pending validation.
        On the coordinator side the diagram highlights the three main outcomes - accept, reject or request clarifications - 
        together with the corresponding status transitions and notifications, matching the feedback loop defined in the requirements.
    }

    {
        The third diagram covers the seriation and assignment process across students, coordinators, companies/professors,
        and the system. 
        It starts with students registering in an edition and ranking proposals, followed by the coordinator triggering
        seriation and the system computing assignments, with different mapping rules for undergraduate internships 
        (one student to zero or one internship) and Master's dissertation (one student to zero or more themes).
        The diagram then shows notifications to students and submitters, interview and confirmation steps on the
        company/professor side, and iterative seriation runs when assignments are not confirmed, or there are still students
        without any assignments.
    }

    {
        Overall, these BPMN models do not attempt to exhaustively represent every alternative or error condition but instead
        provide a high-level, visual view of the core lifecycle processes that drive the platform's requirements and architecture.
        They will be refined in future iterations as the implementation uncovers additional edge cases and 
        as coordinators and other stakeholders provide feedback on real usage. 
    }

\section{Data Model (ERD)}
\label{sec:sol_erd}
\ntindex[Template]{}
    {
        TODO: presentations of the ER diagram (db diagram) e definition of the entities there
    }



\section{Work Plan}
\label{sec:sol_workplan}
\ntindex[Template]{}
    {
        TODO: maybe a cronogram (gantt?) with the development phases, tests, docs writing, risks and mitigations
    }


\iffalse


The third diagram covers the seriation and assignment process across students, coordinators, companies/professors, 
and the system. It starts with students registering in an edition and ranking proposals, followed by the 
coordinator triggering 
seriation and the system computing assignments, with different mapping rules for undergraduate internships 
(one student to zero or one internship) and Master’s dissertations (one student to zero or more themes). 
The diagram then shows notifications to students and submitters, interview and confirmation steps on the 
company/professor side, and iterative seriation runs when assignments are not confirmed, aligning with 
UC‑01, UC‑03, UC‑07, and UC‑08.

Overall, these BPMN models do not attempt to exhaustively represent every alternative or error condition 
but instead provide a high-level, visual view of the core lifecycle processes that drive the platform’s 
requirements and architecture. They will be refined in future iterations as the implementation uncovers additional 
edge cases and as coordinators and other stakeholders provide feedback on real usage.
\fi